% GENERATED FILE. Edit canonical lessons, then run npm run generate:books.
% canonical-source-hash: fnv1a64:b7c5308e3c0ad8e4
% canonical-lessons: AR-C420-istajaba, AR-C420-tajahala, AR-C420-rashshaha, AR-C420-tasawa, AR-C420-mushtarak

\chapter{Responding}
\label{ch:ar-responding}
\hlchaptermodality{420}

\begin{chapteropening}
\textbf{By the end of this chapter:} \emph{I can say that someone responded, ignored, nominated or was equal, and say what is shared.}

Complete the last lesson of chapter 420: I can say that someone responded, ignored, nominated or was equal, and say what is shared.
\end{chapteropening}

\section[\texorpdfstring{\ar{استجاب} (istajāba)}{استجاب (istajāba)}]{\ar{استجاب} (istajāba) — to respond}

\label{lesson:AR-C420-istajaba}



\begin{quote}
Before the new one: say the Arabic for secondary, then the Arabic for comprehensive.
\end{quote}

\subsection*{You'll want to know: \ar{استجاب}}
\textbf{\ar{استجاب}} — \emph{istajāba} — \textquotedblleft{}to respond\textquotedblright{}.

\begin{quote}
\ar{استجاب} \ar{المدير} \ar{لطلب} \ar{العمال}.
\end{quote}

\begin{quote}
\emph{istajāba al-mudīr li-ṭalab al-ʿummāl.}
\end{quote}

\begin{quote}
The manager responded to the workers' request.
\end{quote}

Present: \ar{يستجيب}, \emph{yastajīb}.

It takes \ar{لـ}: \ar{استجاب} \ar{لـ} ….

\begin{cousinweb}
\textbf{\ar{استجاب}} (\emph{istajāba}) is built on the native Arabic root \emph{j-w-b}, \textquotedblleft{}to cut through; to answer\textquotedblright{} (Lane). It is Form X, to give the answer sought: to respond. \ar{جواب}, an answer, is from the same root.
\end{cousinweb}

\begin{grammarlens}[title={What you've built}]
One more word for saying what happens if.
\end{grammarlens}

\subsection*{Your turn}
\begin{itemize}
\raggedright
  \item \emph{Say it:} \emph{istajāba}
  \item \emph{Say it:} \emph{istajāba}, once more
  \item \emph{Say it:} \emph{istajāba al-mudīr li-ṭalab al-ʿummāl.} — The manager responded to the workers' request.
  \item \emph{Recall:} say the Arabic for secondary, then the Arabic for comprehensive, then say \emph{istajāba} again
\end{itemize}

\begin{tcolorbox}[breakable,colback=teal!4,colframe=teal!35!black,title={Before you move on}]
What does \ar{استجاب} mean? (\textquotedblleft{}To respond\textquotedblright{}.) Say it once more.
\end{tcolorbox}
\section[\texorpdfstring{\ar{تجاهل} (tajāhala)}{تجاهل (tajāhala)}]{\ar{تجاهل} (tajāhala) — to ignore}

\label{lesson:AR-C420-tajahala}



\begin{quote}
Before the new one: say the Arabic for comprehensive, then the Arabic for to respond.
\end{quote}

\subsection*{You'll want to know: \ar{تجاهل}}
\textbf{\ar{تجاهل}} — \emph{tajāhala} — \textquotedblleft{}to ignore\textquotedblright{}.

\begin{quote}
\ar{تجاهل} \ar{الولد} \ar{كلام} \ar{أمه}.
\end{quote}

\begin{quote}
\emph{tajāhala al-walad kalām ummih.}
\end{quote}

\begin{quote}
The boy ignored his mother's words.
\end{quote}

Present: \ar{يتجاهل}, \emph{yatajāhal}.

\begin{cousinweb}
\textbf{\ar{تجاهل}} (\emph{tajāhala}) is built on the native Arabic root \emph{j-h-l}, usually traced to the sense \textquotedblleft{}not to know, to be ignorant\textquotedblright{}. It is Form VI, to act as if one did not know: to ignore.
\end{cousinweb}

\begin{grammarlens}[title={What you've built}]
One more word for saying what happens if.
\end{grammarlens}

\subsection*{Your turn}
\begin{itemize}
\raggedright
  \item \emph{Say it:} \emph{tajāhala}
  \item \emph{Say it:} \emph{tajāhala}, once more
  \item \emph{Say it:} \emph{tajāhala al-walad kalām ummih.} — The boy ignored his mother's words.
  \item \emph{Recall:} say the Arabic for comprehensive, then the Arabic for to respond, then say \emph{tajāhala} again
\end{itemize}

\begin{tcolorbox}[breakable,colback=teal!4,colframe=teal!35!black,title={Before you move on}]
What does \ar{تجاهل} mean? (\textquotedblleft{}To ignore\textquotedblright{}.) Say it once more.
\end{tcolorbox}
\section[\texorpdfstring{\ar{رشح} (rashshaḥa)}{رشح (rashshaḥa)}]{\ar{رشح} (rashshaḥa) — to nominate}

\label{lesson:AR-C420-rashshaha}



\begin{quote}
Before the new one: say the Arabic for to respond, then the Arabic for to ignore.
\end{quote}

\subsection*{You'll want to know: \ar{رشح}}
\textbf{\ar{رشح}} — \emph{rashshaḥa} — \textquotedblleft{}to nominate\textquotedblright{}.

\begin{quote}
\ar{رشح} \ar{المدير} \ar{أختي} \ar{لوظيفة} \ar{جديدة}.
\end{quote}

\begin{quote}
\emph{rashshaḥa al-mudīr ukhtī li-waẓīfa jadīda.}
\end{quote}

\begin{quote}
The manager nominated my sister for a new job.
\end{quote}

Present: \ar{يرشح}, \emph{yurashshiḥ}.

It takes \ar{لـ}: \ar{رشحه} \ar{لـ} ….

\begin{cousinweb}
\textbf{\ar{رشح}} (\emph{rashshaḥa}) is built on the native Arabic root \emph{r-sh-ḥ}, \textquotedblleft{}to sweat, to ooze\textquotedblright{} (Lane). It is Form II; Lane gives it for training young ones until they are strong: to put someone forward, to nominate.
\end{cousinweb}

\begin{grammarlens}[title={What you've built}]
One more word for saying what happens if.
\end{grammarlens}

\subsection*{Your turn}
\begin{itemize}
\raggedright
  \item \emph{Say it:} \emph{rashshaḥa}
  \item \emph{Say it:} \emph{rashshaḥa}, once more
  \item \emph{Say it:} \emph{rashshaḥa al-mudīr ukhtī li-waẓīfa jadīda.} — The manager nominated my sister for a new job.
  \item \emph{Recall:} say the Arabic for to respond, then the Arabic for to ignore, then say \emph{rashshaḥa} again
\end{itemize}

\begin{tcolorbox}[breakable,colback=teal!4,colframe=teal!35!black,title={Before you move on}]
What does \ar{رشح} mean? (\textquotedblleft{}To nominate\textquotedblright{}.) Say it once more.
\end{tcolorbox}
\section[\texorpdfstring{\ar{تساوى} (tasāwā)}{تساوى (tasāwā)}]{\ar{تساوى} (tasāwā) — to be equal}

\label{lesson:AR-C420-tasawa}



\begin{quote}
Before the new one: say the Arabic for to ignore, then the Arabic for to nominate.
\end{quote}

\subsection*{You'll want to know: \ar{تساوى}}
\textbf{\ar{تساوى}} — \emph{tasāwā} — \textquotedblleft{}to be equal\textquotedblright{}.

\begin{quote}
\ar{تساوى} \ar{الفريقان} \ar{في} \ar{النقاط}.
\end{quote}

\begin{quote}
\emph{tasāwā al-farīqān fī an-niqāṭ.}
\end{quote}

\begin{quote}
The two teams were level on points.
\end{quote}

Present: \ar{يتساوى}, \emph{yatasāwā}.

\begin{cousinweb}
\textbf{\ar{تساوى}} (\emph{tasāwā}) is built on the native Arabic root \emph{s-w-y}, \textquotedblleft{}to be equal, level\textquotedblright{} (Lane). It is Form VI, to be level with one another: to be equal.
\end{cousinweb}

\begin{grammarlens}[title={What you've built}]
One more word for saying what happens if.
\end{grammarlens}

\subsection*{Your turn}
\begin{itemize}
\raggedright
  \item \emph{Say it:} \emph{tasāwā}
  \item \emph{Say it:} \emph{tasāwā}, once more
  \item \emph{Say it:} \emph{tasāwā al-farīqān fī an-niqāṭ.} — The two teams were level on points.
  \item \emph{Recall:} say the Arabic for to ignore, then the Arabic for to nominate, then say \emph{tasāwā} again
\end{itemize}

\begin{tcolorbox}[breakable,colback=teal!4,colframe=teal!35!black,title={Before you move on}]
What does \ar{تساوى} mean? (\textquotedblleft{}To be equal\textquotedblright{}.) Say it once more.
\end{tcolorbox}
\section[\texorpdfstring{\ar{مشترك} (mushtarak)}{مشترك (mushtarak)}]{\ar{مشترك} (mushtarak) — shared}

\label{lesson:AR-C420-mushtarak}



\begin{quote}
Before the new one: say the Arabic for to nominate, then the Arabic for to be equal.
\end{quote}

\subsection*{You'll want to know: \ar{مشترك}}
\textbf{\ar{مشترك}} — \emph{mushtarak} — \textquotedblleft{}shared\textquotedblright{}.

\begin{quote}
\ar{عندنا} \ar{أصدقاء} \ar{مشتركون}.
\end{quote}

\begin{quote}
\emph{ʿindanā aṣdiqāʾ mushtarakūn.}
\end{quote}

\begin{quote}
We have friends in common.
\end{quote}

\begin{cousinweb}
\textbf{\ar{مشترك}} (\emph{mushtarak}) is built on the native Arabic root \emph{sh-r-k}, \textquotedblleft{}to share, to be a partner\textquotedblright{} (Lane). \emph{Mushtarak} is the Form VIII passive participle, what is shared: common, joint.
\end{cousinweb}

\begin{grammarlens}[title={What you've built}]
One more word for saying what happens if.
\end{grammarlens}

\subsection*{Your turn}
\begin{itemize}
\raggedright
  \item \emph{Say it:} \emph{mushtarak}
  \item \emph{Say it:} \emph{mushtarak}, once more
  \item \emph{Say it:} \emph{ʿindanā aṣdiqāʾ mushtarakūn.} — We have friends in common.
  \item \emph{Recall:} say the Arabic for to nominate, then the Arabic for to be equal, then say \emph{mushtarak} again
\end{itemize}

\begin{tcolorbox}[breakable,colback=teal!4,colframe=teal!35!black,title={Before you move on}]
What does \ar{مشترك} mean? (\textquotedblleft{}Shared\textquotedblright{}.) Say it once more.
\end{tcolorbox}
